\documentclass{article} % Default font size and left-justified equations
% \documentclass{book}

\usepackage[
  margin=0.7in,
  % includefoot,
  % footskip=30pt,
]{geometry}
\usepackage{listings}
\usepackage{color}
\usepackage{natbib}
\usepackage{enumitem}
\usepackage{amsmath,amsfonts}
\usepackage{graphicx}
\usepackage{framed}
\usepackage{subcaption}
\usepackage[T1]{fontenc}%
\usepackage[utf8]{inputenc}%
\usepackage{mathrsfs}%
\usepackage{amssymb}%
\usepackage{amsthm}%
\usepackage{graphicx}
\usepackage{hyperref}
\usepackage{environ}
\usepackage{framed}
\usepackage{mdframed} 
\usepackage{wrapfig} 
\usepackage{booktabs}
\usepackage{tabularx} 
\usepackage{array}
\usepackage[font=small]{caption} 
\usepackage{xspace}  
\usepackage[osf,sc]{mathpazo}
\usepackage{pbox}
\usepackage{listings}
\usepackage{tikz}
\usepackage{bm}
\usepackage[strict]{changepage}
\usepackage{mleftright}
\usepackage{appendix}
\usepackage{wrapfig}
\usepackage[all]{xy} 
\setcounter{tocdepth}{3}
\setcounter{secnumdepth}{3}
% \usepackage[ruled,vlined]{algorithm2e}
\usepackage{algorithm}
\usepackage{algpseudocode}
\usepackage{framed}
\usepackage{wrapfig}
\usepackage{pythonhighlight}
\usepackage{tikz}
\usetikzlibrary{arrows}
\usetikzlibrary{arrows.meta}
\usetikzlibrary{shapes.geometric}
\usetikzlibrary{positioning, arrows, automata, calc}
\usepackage{transparent}
\usepackage[many]{tcolorbox}
\usepackage{tikz}
\usetikzlibrary{shapes.geometric}
\usetikzlibrary{positioning, arrows, automata, calc}

\newtcolorbox[]{your_solution}[1][]{
    % breakable,
    enhanced,
    nobeforeafter,
    colback=white,
    title=Your Answer,
    sidebyside align=top,
    box align=top,
    #1
}
\newcommand{\mybox}[1]{
\noindent
\fbox{\parbox{0.955\textwidth}{%
\noindent\texttt{#1}
}
}
}

\lstset{
    % backgroundcolor=\color{lightgray},   % choose the background color
    basicstyle=\ttfamily\small,          % font style
    frame=single,                        % add a frame around the code
    breaklines=true,                     % automatic line breaking
    captionpos=b,                        % position of the caption
    numbers=left,                        % where to put the line-numbers
    numberstyle=\tiny\color{gray},       % the style of the line numbers
    keywordstyle=\color{blue},           % keyword style
    commentstyle=\color{green},          % comment style
    stringstyle=\color{red},             % string literal style
}

\newcommand{\filler}{ . . . . . }
\newcommand{\choice}{\hspace{0.5cm}$\square$}
\newcommand{\identity}{\mathbf{I}} 
\newcommand{\paran}[1]{\left( #1 \right)}

%% PLEASE USE THESE MACROS WHEN WRITING THE SOLUTIONS 
\long\def\solution#1{\par\noindent{\color{purple}Answer: \em #1}\par} % show
%\newcommand{\solution}[2]{#2} % hide

\newcommand{\extracredit}{{\color{purple}\textbf{Extra Credit:}}}

\newcommand{\additionalNotes}{
\noindent\textbf{How to hand in your written work and code:} via Gradescope.  \\ 


% \noindent\textbf{Collaboration:} Make certain that you understand the course collaboration policy, described on the course website. 
% You may discuss the homework to understand the problems and the mathematics behind the various learning algorithms, but you are \textcolor{red}{\textbf{not allowed to share solutions with any other students. You must write the solutions \textbf{individually}}}. \\ 

% \noindent\textbf{Typesetting:} 
% We strongly recommend typesetting your homework, especially if you have sloppy handwriting! :)  
% % We recommend using LaTeX, but you are welcome to use whatever program and/or markup language you like. 
% We will provide a LaTeX template for homework solutions.
}

\newcommand{\additionalNotesTeamUpThree}{
\noindent\textbf{How to hand in your written work:} via Gradescope.  \\ 


\noindent\textbf{Collaboration:} 
You may discuss the homework with your classmates, if that helps you understand the problems and various ideas in the class, but you are \textcolor{red}{\textbf{not allowed to share problem solutions with any other students. You must write the solutions \textbf{individually}}. 
% You're not allowed to use AI for drafting responses.
The use of AI (e.g., Github CoPilot) is permitted but \textbf{must be documented}. 
} 
% Using AI for generating code (e.g., Github CoPilot) is permitted. 
\\ 

\noindent\textbf{Typesetting:} 
We strongly recommend typesetting your homework, especially if you have sloppy handwriting! :)  
% We recommend using LaTeX, but you are welcome to use whatever program and/or markup language you like. 
We will provide a LaTeX template for homework solutions.
Please be as succinct as possible in your answers; overly long responses will be penalized.

}

% \newcommand{\additionalNotesTeamUpThree}{
% \noindent\textbf{How to hand in your written work:} via Gradescope.  \\ 

\title{ \Large (CS 601.768) Language Model Agents }
\date{For homework deadline and collaboration policy, check the syllabus \\
\normalsize Name: Dev Vyas\_\_\_\_\   \\ 
% Collaborators, if any: \_\_\_\_\_\_\_\_\_\_\_\_\_\_\_\_\_\_\_\_\_\_\_\_\_\_\_\_\_\_\_\_\_\_\_\_\_\_\_\_\_\_\_\_\_\_\_\_\_\_\_\_ \\ 
Sources used for your homework, if any: \_\_\_\_\_\_\_\_\_\_\_\_\_\_\_\_\_\_\_\_\_\_\_\_\_\_\_\_\_\_\_\_\_\_\_\_\_\_\_\ 
}
\newcommand{\todo}{\textcolor{blue}{\textbf{TODOs}}}

